\section{Conclusion}

In this work, we investigated how LLMs represent the computationally opaque vigesimal counting system of \standardfrench compared to the transparent decimal system of \swissfrench. We demonstrated that modern SOTA models have successfully overcome the ``vigesimal tax'' in zero-shot decoding, translating both dialects into digits with perfect accuracy. However, this surface-level success masks deeper reasoning misalignments. When forced to compare the internal representations of these numbers across dialects, models exhibit significant error rates and novel biases, such as the ``Standard Overestimation'' bias, where linguistic complexity is conflated with magnitude. Ultimately, our findings reveal that while models can translate complex structural forms, mapping these linguistic structures to a consistent abstract number line remains an unsolved challenge. Future work should probe the residual streams of these models to isolate the specific circuits responsible for vigesimal versus decimal processing, and explore whether fine-tuning on transparent decimal variants can generally improve multilingual mathematical reasoning.